\section{An affine invariant with explicit constants}\label{s:affine-invariant}

The radius-weighted norm packages every test of one dimension into a
single number. A closely related formulation propagates an explicit
bound on every affine subsphere at once. It uses a different moment
order and an even radius exponent, so its inverse-radius calculation
is an exact finite beta product. We retain its numerical constants.

In this section \(m\) is the block length and \(k\) is the number of
blocks. Set
\begin{equation}\label{s:affine-parameters}
 m=\lfloor d/16\rfloor,\qquad p=\lfloor d/8\rfloor,\qquad
 a=2\lfloor d/32\rfloor,\qquad n_0=\lceil d/2\rceil,\qquad
 C_*=1000^{32}.
\end{equation}

\begin{proposition}[All-affine suffix bound]\label{s:affine-proposition}
For all sufficiently large \(d\), consider a Borel learner padded to
\(k\) blocks of \(m\) exact Gaussian samples, of width at most \(W\),
and with target-independent unit output laws at the last layer.
The randomized experiment, shared seed, initialization, and fresh
randomness satisfy the joint measurability, independence, and
information restrictions of Definition~\ref{s:model}. Here
\(0<\epsilon\le1/10\), and success means angular error at most
\(\epsilon\).
Assume
\begin{equation}\label{s:affine-budget}
 n_0+2k\le d,\qquad W^{1/p}\le2^d.
\end{equation}
For each fixed shared seed whose rules satisfy these assumptions, and
for a state \(v\) at layer \((k-\ell)m\), let \(g_v(s)\) be the
success probability of the suffix started in that state. For
\(0\le\ell\le k\), every affine subsphere \eqref{s:affine-sphere}
with \(\dim H=n\ge n_0+2\ell\) satisfies
\begin{equation}\label{s:affine-invariant-bound}
 \int g_v(z+r\xi)\,d\sigma_H(\xi)
 \le (R_\ell/r)^a,\qquad R_\ell=16C_*^\ell\epsilon.
\end{equation}
The assertion is uniform in the fixed seed and includes states
unreachable from the initial layer.
\end{proposition}
\begin{proof}
At \(\ell=0\), fix a terminal output. If \(16\epsilon/r\ge1\),
total mass proves the bound. Otherwise the successful directions lie
in a chordal cap of radius \(2\epsilon/r\). Its polar angle is at most
\(4\epsilon/r\), so the polar calculation in Lemma~\ref{s:caps}
bounds its mass by \((16\epsilon/r)^{n-1}\).
Since \(n-1\ge a\), this is at most \((16\epsilon/r)^a\).
Average over the output law.

Suppose the assertion holds at \(\ell\). Consider a preceding block
on an \(n\)-dimensional test with \(n\ge n_0+2(\ell+1)\).
Lemma~\ref{s:residual}, with \(b=m\), gives a beta fraction \(t\),
a direction \(e\), and a uniform residual \(K=n-m\) dimensional
space \(L\) inside \(E=H\cap e^\perp\), of dimension \(N=n-1\).
Put
\[
 z'=z+r\sqrt{1-t}\,e,\qquad \rho=r\sqrt t,\qquad
 f_u(y)=g_u(z'+\rho y),\qquad F_u(L)=\int f_u\,d\sigma_L.
\]
By Lemma~\ref{s:state-maximum}, the desired mean is at most
\begin{equation}\label{s:affine-state-max}
 \E_{t,e}\E_L\max_u F_u(L).
\end{equation}
The maximum is over all next boundary states and has already removed
their possible dependence on the rest of the exact block data.

Fix \(t,e,u\) and write \(\Lambda=(R_\ell/\rho)^a\). The test
\(z'+\rho S(E)\) is admissible and has dimension \(n-1\), so the
induction hypothesis gives \(\int f_u\,d\sigma_E\le\Lambda\).
For a cap about \(y_0\in S(E)\), condition a uniform direction on
its polar angle \(\theta\). The corresponding target is uniform on
the affine subsphere with center \(z'+\rho\cos\theta\,y_0\), radius
\(\rho\sin\theta\), and space \(E\cap y_0^\perp\) of dimension
\(n-2\ge n_0+2\ell\). Its average is at most
\(\Lambda(\sin\theta)^{-a}\).

For \(0<\delta\le1\), the chordal cap has polar angle at most
\(2\delta\). Integrate the latitude bound against the normalized
\(\sin^{N-2}\theta\) density. The exponent \(N-2-a\) is
nonnegative, and the original normalizing integral is at least
\(2^{-(N-2)}\). Using \(\sin\theta\le\theta\) gives
\begin{equation}\label{s:affine-cap}
 \int_{\norm{y-y_0}\le\delta}f_u\,d\sigma_E
 \le4^d\Lambda\delta^{N-1-a}.
\end{equation}
For \(\delta>1\), total mass proves the same bound.

Apply Lemma~\ref{s:cap-moment} with \(A_0=1,A_1=4\),
moment \(p\), and singular exponent \(N-K=m-1\). We have
\(1\le p<K\le N\), \(0\le a<N-1\), and, for \(0\le j<p\),
\[
 N-1-a-j-(m-1)
 \ge n-a-p-m
 \ge d/4+2(\ell+1)\ge1.
\]
Indeed \(n\ge n_0+2(\ell+1)\ge d/2+2(\ell+1)\), while
\(a\le d/16\), \(p\le d/8\), and \(m\le d/16\).
The tube coefficient in \eqref{s:tube} is \(36^d\), and
\eqref{s:inverse-distance} gives \(144^d\Lambda\). Consequently
\[
 \E_L F_u(L)^p\le(144^d\Lambda)^p,\qquad
 \E_L\max_uF_u(L)\le W^{1/p}144^d\Lambda\le288^d\Lambda.
\]

The remaining radius fraction has
\(t\sim\Beta((n-m)/2,m/2)\). Here \(q=a/2=\lfloor d/32\rfloor\)
is an integer. Its denominator shapes satisfy
\[
 (n-m)/2-q\ge3d/16>0.
\]
Since \(n/2\le d/2\), every ratio in the finite product is at most
\((d/2)/(3d/16)=8/3<3\). Lemma~\ref{s:beta-identities} gives
\begin{equation}\label{s:affine-beta}
 \E t^{-a/2}
 =\prod_{j=1}^{q}\frac{n/2-j}{(n-m)/2-j}\le3^d.
\end{equation}
Substitution into \eqref{s:affine-state-max} yields
\(864^d(R_\ell/r)^a\). For large \(d\), \(a\ge d/32\), hence
\(C_*^a\ge1000^d\ge864^d\). This proves the induction step with
\(R_{\ell+1}=C_*R_\ell\).
\end{proof}

To apply the proposition, let \(h=\log_2(1/\epsilon)\), assume
\(M(d)=o(d^2)\), and suppose uniform-prior success is at least \(2/3\).
The dimensions with \(T>dh\) already have the desired rate. On those
with \(T\le dh\), apply Lemma~\ref{s:short-run} with
\(c=1\) and \(\alpha=2/3\). It gives \(h=o(d)\),
\(T=o(d^2)\), and padded log-width \(o(d^2)\).
Pad to \(k=\lceil T/m\rceil=o(d)\) blocks with
\(W=(T+2)2^M\). Then \eqref{s:affine-budget} holds for all
sufficiently large such dimensions: \(n_0+2k\le d\), and
\(\log_2 W=o(d^2)\) with \(p\ge d/16\) eventually gives
\(W^{1/p}\le2^d\).

At the initial sphere \(z=0,r=1,H=\R^d\), average the bound over
initialization conditional on the shared seed, and then over that
seed. Since \(a\ge1\),
\begin{equation}\label{s:affine-endgame}
 \frac23\le(16C_*^k\epsilon)^a,\qquad
 h\le k\log_2C_*+\log_2 24
 \le(T/m+1)\log_2C_*+\log_2 24.
\end{equation}
For \(h\ge2(\log_2C_*+\log_2 24)\), this gives
\[
 T\ge \frac{mh}{2\log_2C_*},\qquad m\ge d/32
\]
eventually. For smaller \(h\), \eqref{s:quarter-endpoint} gives
the same \(d h\) rate with an absolute constant. This completes the
all-affine proof, including its terminal dimension and
bounded-precision cases.
